% Options for packages loaded elsewhere
\PassOptionsToPackage{unicode}{hyperref}
\PassOptionsToPackage{hyphens}{url}
\documentclass[
]{article}
\usepackage{xcolor}
\usepackage{amsmath,amssymb}
\setcounter{secnumdepth}{-\maxdimen} % remove section numbering
\usepackage{iftex}
\ifPDFTeX
  \usepackage[T1]{fontenc}
  \usepackage[utf8]{inputenc}
  \usepackage{textcomp} % provide euro and other symbols
\else % if luatex or xetex
  \usepackage{unicode-math} % this also loads fontspec
  \defaultfontfeatures{Scale=MatchLowercase}
  \defaultfontfeatures[\rmfamily]{Ligatures=TeX,Scale=1}
\fi
\usepackage{lmodern}
\ifPDFTeX\else
  % xetex/luatex font selection
\fi
% Use upquote if available, for straight quotes in verbatim environments
\IfFileExists{upquote.sty}{\usepackage{upquote}}{}
\IfFileExists{microtype.sty}{% use microtype if available
  \usepackage[]{microtype}
  \UseMicrotypeSet[protrusion]{basicmath} % disable protrusion for tt fonts
}{}
\makeatletter
\@ifundefined{KOMAClassName}{% if non-KOMA class
  \IfFileExists{parskip.sty}{%
    \usepackage{parskip}
  }{% else
    \setlength{\parindent}{0pt}
    \setlength{\parskip}{6pt plus 2pt minus 1pt}}
}{% if KOMA class
  \KOMAoptions{parskip=half}}
\makeatother
\usepackage{longtable,booktabs,array}
\usepackage{caption}
\captionsetup[table]{skip=6pt}
\newcounter{none} % for unnumbered tables
\usepackage{calc} % for calculating minipage widths
% Correct order of tables after \paragraph or \subparagraph
\usepackage{etoolbox}
\makeatletter
\patchcmd\longtable{\par}{\if@noskipsec\mbox{}\fi\par}{}{}
\makeatother
\setlength{\emergencystretch}{3em} % prevent overfull lines
\providecommand{\tightlist}{%
  \setlength{\itemsep}{0pt}\setlength{\parskip}{0pt}}
\usepackage{bookmark}
\IfFileExists{xurl.sty}{\usepackage{xurl}}{} % add URL line breaks if available
\urlstyle{same}
% fallback for those not using the hyperref driver hyperxmp:
\makeatletter
\@ifundefined{xmpquote}{\newcommand{\xmpquote}[1]{#1}}{}
\makeatother
\hypersetup{
  hidelinks,
  pdfcreator={LaTeX via pandoc}}

\author{}
\date{}

\begin{document}

{
\setcounter{tocdepth}{2}
\tableofcontents
}
\section{Programarea pe blocuri în WebCoder --- Blockly ȘI
Scratch}\label{programarea-pe-blocuri-uxeen-webcoder--blockly-ux219i-scratch}

\begin{quote}
Directivă owner (2026-10-30): platforma livrează \textbf{ambele} motoare
de blocuri --- Blockly \emph{și} Scratch --- nu unul sau altul. Blockly
e puntea structurată spre codul tipizat; Scratch e stratul de
familiaritate (pisica Scratch, scena, say/answer) pentru cei mai mici
învățăcei. Acest document este ghidul complet al amândurora:
arhitectură, integrarea cu evaluatorul, authoring, învățare, predare,
testare, securitate, troubleshooting și exemple lucrate. E scris astfel
încât un contributor nou să poată extinde sistemul fără să citească mai
întâi codul. Acesta este rezumatul executiv al ghidului; ghidul tehnic
complet, în engleză, este
\href{../guide-en/BLOCK-PROGRAMMING.md}{\texttt{docs/guide-en/BLOCK-PROGRAMMING.md}}.
\end{quote}

\begin{center}\rule{0.5\linewidth}{0.5pt}\end{center}

\subsection{1. Ce s-a livrat, pe scurt}\label{1-ce-s-a-livrat-pe-scurt}

O problemă poate fi o \textbf{problemă cu blocuri}. Creatorul alege un
\emph{motor} de blocuri per problemă:

\begin{itemize}
\tightlist
\item
  \textbf{Blockly}
  (\texttt{block\_config.engine:\ \textquotesingle{}blockly\textquotesingle{}},
  implicit) --- editor Blockly cu un toolbox curat de WebCoder
  (input/control/logic/math/text/lists/variables), care
  \textbf{generează cod} (implicit Python) trimis evaluatorului exact ca
  o submissions obișnuită --- verdict real, AC/WA, tot pipeline-ul;
\item
  \textbf{Scratch}
  (\texttt{block\_config.engine:\ \textquotesingle{}scratch\textquotesingle{}})
  --- proiect Scratch (SB3) încărcat de elev; un \textbf{transpiler} îl
  transformă în Python înainte de evaluare (transpile doar pentru
  evaluare --- proiectul rămâne sursa de adevăr).
\end{itemize}

Ambele motoare folosesc aceeași structură \texttt{block\_config} pe
problemă, același editor parțial cu regiuni blocate (PartialCodeEditor)
și aceleași limite ale evaluatorului.

\subsection{2. De ce ambele}\label{2-de-ce-ambele}

\begin{itemize}
\tightlist
\item
  \textbf{Blockly} pregătește structura tipizată: variabile, condiții,
  bucle --- pasul natural spre Python/C++.
\item
  \textbf{Scratch} aduce familiaritatea: elevii de gimnaziu știu deja
  scena și pisica; teama dispare.
\item
  Profesorul alege per problemă; curriculumul Scratch→Python
  (cursul-punte din \texttt{/courses}) folosește ambele.
\end{itemize}

\subsection{3. Fluxul elevului}\label{3-fluxul-elevului}

\begin{enumerate}
\def\labelenumi{\arabic{enumi}.}
\tightlist
\item
  Deschide o problemă cu blocuri → editorul de blocuri înlocuiește
  editorul de text (picker-ul de limbaj dispare).
\item
  Construiește programul prin tragere de blocuri; toolbox-ul e curat per
  problemă (creatorul poate restriciona categoriile).
\item
  \textbf{Trimite} --- codul generat (Blockly) sau transpilat (Scratch)
  merge prin pipeline-ul normal de evaluare; UI-ul face polling până la
  verdictul final (specificația real-stack e2e acoperă integral acest
  flux).
\item
  Vederea verdictului arată același tabel per test ca la problemele
  obișnuite.
\end{enumerate}

\subsection{4. Fluxul
profesorului/creatorului}\label{4-fluxul-profesoruluicreatorului}

\begin{itemize}
\tightlist
\item
  La authoring, tabul de blocuri din wizard setează motorul, generatorul
  (limbajul-țintă pentru Blockly) și toolbox-ul;
\item
  problemele cu cod parțial (regiuni blocate + regiuni de completat)
  funcționează în ambele motoare --- "completează bucla lipsă";
\item
  suita de teste se definește identic (input/expected per test).
\end{itemize}

\subsection{5. Detalii tehnice
(rezumat)}\label{5-detalii-tehnice-rezumat}

\begin{itemize}
\tightlist
\item
  \texttt{Problem.block\_config} (JSON): \texttt{engine},
  \texttt{generator}, \texttt{toolbox} (lista de categorii), opțional
  \texttt{starter} (programul de start) și \texttt{locked\_regions}
  pentru cod parțial;
\item
  generarea Blockly: blocuri → cod sursă pe client, trimis ca
  \texttt{code} obișnuit (limbajul = generatorul);
\item
  transpilarea Scratch: \texttt{problems/scratch\_transpiler.py}
  (server-side, deterministic);
\item
  securitate: nimic din config nu execută cod pe server; tot ce rulează
  trece prin sandbox-ul normal al evaluatorului.
\end{itemize}

Verificați ghidul EN pentru arhitectura completă, exemplele lucrate,
matricea de testare și troubleshooting-ul (571 de linii de detail).

\subsection{6. Testare}\label{6-testare}

\begin{itemize}
\tightlist
\item
  unit: vitest pe generatoare și pe transpiler;
\item
  e2e mock: payload-ul submission-ului generat din blocuri;
\item
  e2e real-stack (\texttt{E2E\_API\_BASE}): round-trip complet cu
  verdicte evaluate prin polling --- problema demo „E2E Blockly Real
  Last Digit".
\end{itemize}

\subsection{7. Probleme frecvente}\label{7-probleme-frecvente}

{\def\LTcaptype{none} % do not increment counter
\begin{longtable}[]{@{}ll@{}}
\toprule\noalign{}
Simptom & Cauză probabilă \\
\midrule\noalign{}
\endhead
\bottomrule\noalign{}
\endlastfoot
Editorul de blocuri nu apare & problema nu are \texttt{block\_config};
verifică tabul de blocuri din wizard \\
Codul generat nu compilează & generatorul nu corespunde limbajului
problemei (setează \texttt{generator}) \\
Proiectul Scratch e respins & SB3 invalid sau blocuri netranspilabile;
vezi lista de susținere din ghidul EN \\
\end{longtable}
}

\end{document}
